
\section{Coheight--Krull dim bridge for scheme points}
\label{mk-chap:Albanese_CoheightBridge}

\section{Coheight is preserved by open immersions}
\label{mk-sec:coheight_topology}

\begin{lemma}[Coheight on opens equals coheight in the ambient space]\label{mk-lem:coheight_eq_of_isOpenEmbedding}
  \dcref{custom}

  Let \(X\) be a topological space carrying the specialisation preorder, and
  let \(U \subseteq X\) be an open subset. For any \(z \in U\),
  \[
    \operatorname{coheight}_{X}(z) \;=\; \operatorname{coheight}_{U}\bigl(\langle z, \, z \in U\rangle\bigr),
  \]
  where on the right \(U\) carries the subspace topology and \(\langle z, \,
  z \in U\rangle\) denotes \(z\) viewed as a point of \(U\).
\end{lemma}

\begin{proof}

Both sides are by definition the supremum, over natural numbers \(n\), of the
  set of \(n\) for which there exists a strictly increasing chain
  \(z = a_0 \prec a_1 \prec \cdots \prec a_n\) in the relevant
  specialisation preorder (with \(\prec\) denoting strict specialisation). It
  therefore suffices to give an order-preserving bijection between the
  strict-generisation chains starting at \(z\) in \(X\) and those starting at the
  corresponding point of \(U\).

  \emph{Chains in \(U\) embed into \(X\).} The inclusion \(U \hookrightarrow X\) is
  continuous, so it preserves specialisation: if \(y \rightsquigarrow x\) in \(U\)
  then \(y \rightsquigarrow x\) in \(X\), and distinct points of \(U\) remain
  distinct in \(X\). Hence every strict-generisation chain in \(U\) starting at
  \(z\) lifts to one in \(X\) of the same length.

  \emph{Chains in \(X\) starting at \(z\) all lie in \(U\).} Conversely, suppose
  \(z = a_0 \prec a_1 \prec \cdots \prec a_n\) is a strict-generisation chain
  in \(X\) with \(a_0 = z \in U\). For each \(i \geq 0\), the relation \(a_i \succeq
  z\) (in the specialisation preorder) is \(z \rightsquigarrow a_i\), i.e.\
  \(a_i\) is in the closure of \(\{z\}\) \emph{viewed from above}; in our
  convention the strict order \(a \prec b\) records \(b \rightsquigarrow a\)
  (the generic point is the strict maximum), so \(a_i \succeq z\) in \(X\) means
  \(a_i \rightsquigarrow z\) in \(X\). By the Mathlib lemma
  \texttt{Topology.Inseparable.Specializes.mem\_open} (any specialisation
  \(a \rightsquigarrow z\) with \(z\) in an open set \(U\) forces \(a \in U\), since
  \(U\) is open and contains \(z\), hence contains every point whose closure
  contains \(z\) --- equivalently every \(a\) with \(a \rightsquigarrow z\)),
  every \(a_i\) lies in \(U\). Thus the chain descends to a strict-generisation
  chain \(\langle z\rangle = \langle a_0\rangle \prec \langle a_1\rangle \prec
  \cdots \prec \langle a_n\rangle\) in \(U\) of the same length.

  These two operations are mutually inverse on the set of strict-generisation
  chains starting at \(z\) (resp.\ at \(\langle z, \, z \in U\rangle\)) and
  preserve length, so the suprema agree.
\end{proof}

\section{Coheight on Spec equals height in PrimeSpectrum}
\label{mk-sec:coheight_spec_duality}

\begin{lemma}[Coheight on \(\Spec R\) equals height in \(\operatorname{PrimeSpectrum} R\)]\label{mk-lem:coheight_spec_eq_height_primeSpectrum}
  \dcref{custom}

  Let \(R\) be a commutative ring and let \(p \in \Spec R\) (viewed as a point of
  the scheme spectrum, with the specialisation preorder). Then
  \[
    \operatorname{coheight}_{\Spec R}(p) \;=\; \operatorname{height}_{\operatorname{PrimeSpectrum}(R)}(p),
  \]
  where the right-hand side is the supremum of the lengths of strict
  prime-ideal chains \(\mathfrak{p}_0 \subsetneq \cdots \subsetneq \mathfrak{p}\)
  ending at the prime \(\mathfrak{p}\) corresponding to \(p\).
\end{lemma}

\begin{proof}

The specialisation preorder on \(\Spec R\) is the dual of the natural
  inclusion preorder on \(\operatorname{PrimeSpectrum}(R)\). Concretely, for
  \(p, q \in \Spec R\) with underlying primes \(\mathfrak{p}, \mathfrak{q}\),
  the relation \(p \leq q\) in the specialisation preorder (which is
  \(q \rightsquigarrow p\) in topological notation) corresponds to the prime
  containment \(\mathfrak{q} \subseteq \mathfrak{p}\). This is the content of
  the Mathlib lemma \texttt{spec\_le\_iff} (located in
  \texttt{Mathlib.AlgebraicGeometry.AffineSpace}, lines 438--447 at commit
  \texttt{b80f227}, with the explicit remark that the preorder on \(\Spec R\)
  equals the order-dual of the preorder on \(\operatorname{PrimeSpectrum} R\)).

  Under this order anti-isomorphism, coheight of \(p\) in \(\Spec R\)
  (= supremum of strict-generisation chain lengths starting at \(p\),
  i.e.\ ascending chains \(p \prec a_1 \prec \cdots \prec a_n\) in the
  specialisation preorder) translates to height of \(\mathfrak{p}\) in
  \(\operatorname{PrimeSpectrum}(R)\) (= supremum of strict descending prime-ideal
  chains \(\mathfrak{p} \supsetneq \mathfrak{q}_1 \supsetneq \cdots \supsetneq
  \mathfrak{q}_n\), equivalently strict ascending chains
  \(\mathfrak{q}_n \subsetneq \cdots \subsetneq \mathfrak{q}_1 \subsetneq
  \mathfrak{p}\)). The transport is mechanised by the Mathlib lemmas
  \texttt{Order.coheight\_orderIso} (coheight is invariant under order
  isomorphisms) and \texttt{Order.height\_toDual} (height in the order-dual
  equals coheight in the original order), applied to the order-dual
  identification \(\Spec R \simeq_o (\operatorname{PrimeSpectrum} R)^{\operatorname{op}}\)
  furnished by \texttt{spec\_le\_iff}.
\end{proof}

\section{The coheight \(=\) Krull-dim-of-stalk bridge}
\label{mk-sec:coheight_bridge}

\begin{theorem}[Coheight-to-Krull-dim bridge for a scheme point]\label{mk-thm:ringKrullDim_stalk_eq_coheight}
  \dcref{custom}

  \uses{mk-lem:coheight_eq_of_isOpenEmbedding, mk-lem:coheight_spec_eq_height_primeSpectrum}

  Let \(X\) be a scheme and let \(z \in X\) be any point. Then
  \[
    \operatorname{ringKrullDim}(\mathcal{O}_{X,z}) \;=\; \operatorname{coheight}_{X}(z),
  \]
  i.e.\ the Krull dimension of the stalk at \(z\) equals the coheight of \(z\) in
  the underlying topological space of \(X\) (with respect to the specialisation
  preorder).
\end{theorem}

\begin{proof}

  The proof proceeds in five steps, assembling Mathlib API rather than
  arguing from a single textbook reference.

  \emph{Step 1: pick an affine open neighbourhood of \(z\).} By
  \texttt{AlgebraicGeometry.exists\_isAffineOpen\_mem\_and\_subset}
  applied to the open set \(X\) itself, there exists an affine open
  \(U \subseteq X\) with \(z \in U\). Write \(\langle z, \, z \in U \rangle \in U\)
  for \(z\) viewed as a point of the subscheme \(U\).

  \emph{Step 2: name the prime corresponding to \(z\) inside \(U\).} Since
  \(U \cong \Spec \Gamma(U, \mathcal{O}_X)\) via the affine equivalence
  \(U \simeq \Spec R\) (with \(R := \Gamma(U, \mathcal{O}_X)\)), the point
  \(\langle z, \, z \in U\rangle\) corresponds to a prime ideal \(\mathfrak{p}
  \subseteq R\). In Mathlib this prime is named
  \texttt{IsAffineOpen.primeIdealOf}\,\(\langle z, \, z \in U\rangle\).

  \emph{Step 3: the stalk is the localisation of \(\Gamma(U)\) at
  \(\mathfrak{p}\).} The classical affine-stalk identification says
  \(\mathcal{O}_{X,z} \cong R_{\mathfrak{p}}\). In Mathlib this is the
  instance \texttt{IsAffineOpen.isLocalization\_stalk}, which witnesses
  \(\mathcal{O}_{X,z}\) as a localisation of \(R\) at \(\mathfrak{p}\) in the
  \texttt{IsLocalization.AtPrime} sense.

  \emph{Step 4: Krull dim of the localisation equals the height of
  \(\mathfrak{p}\).} The Mathlib lemma
  \texttt{IsLocalization.AtPrime.ringKrullDim\_eq\_height} gives
  \[
    \operatorname{ringKrullDim}(\mathcal{O}_{X,z}) \;=\;
    \operatorname{ringKrullDim}(R_{\mathfrak{p}}) \;=\;
    \operatorname{height}(\mathfrak{p}),
  \]
  where \(\operatorname{height}(\mathfrak{p})\) is the prime-ideal height in
  \(\operatorname{PrimeSpectrum}(R)\). By
  \texttt{Ideal.height\_eq\_primeHeight} this matches the
  \(\operatorname{height}_{\operatorname{PrimeSpectrum}(R)}\) used in
  \ref{mk-lem:coheight_spec_eq_height_primeSpectrum}.

  \emph{Step 5: lift back from \(U\) to \(X\).} By
  \ref{mk-lem:coheight_spec_eq_height_primeSpectrum} applied to \(R\) and the
  affine identification \(U \simeq \Spec R\),
  \[
    \operatorname{height}_{\operatorname{PrimeSpectrum}(R)}(\mathfrak{p}) \;=\;
    \operatorname{coheight}_{\Spec R}(\langle\text{image of } z\rangle) \;=\;
    \operatorname{coheight}_{U}\bigl(\langle z, \, z \in U\rangle\bigr).
  \]
  Then \ref{mk-lem:coheight_eq_of_isOpenEmbedding} applied to the open
  \(U \subseteq X\) promotes the right-hand side back to
  \(\operatorname{coheight}_{X}(z)\). Concatenating the equalities from
  Steps~4--5 gives the conclusion.
\end{proof}

\section{Codim-1 specialisation: synthesising \texorpdfstring{\(\operatorname{KrullDimLE} 1\)}{KrullDimLE 1}}
\label{mk-sec:coheight_le_one_instance}

\begin{lemma}[Coheight \(= 1\) forces stalk \(\operatorname{Krull\,dim} \leq 1\)]\label{mk-lem:ringKrullDimLE_of_coheight_eq_one}
  \dcref{custom}

  \uses{mk-thm:ringKrullDim_stalk_eq_coheight}

  Let \(X\) be a scheme and let \(z \in X\) be a point with
  \(\operatorname{coheight}_{X}(z) = 1\). Then the stalk \(\mathcal{O}_{X,z}\)
  has Krull dimension at most \(1\), equivalently \(\operatorname{Ring.KrullDimLE}\,1\,
  \mathcal{O}_{X,z}\) holds.
\end{lemma}

\begin{proof}

  By \ref{mk-thm:ringKrullDim_stalk_eq_coheight},
  \(\operatorname{ringKrullDim}(\mathcal{O}_{X,z}) = \operatorname{coheight}_{X}(z) = 1\).
  The Mathlib predicate \(\operatorname{Ring.KrullDimLE}\,1\,R\) unfolds to
  \(\operatorname{ringKrullDim}(R) \leq 1\), so the inequality \(1 \leq 1\)
  discharges the goal.
\end{proof}
